\documentclass[leqno]{article}
% Shared mark scheme preamble for resources.danieljsmith.org: the pack
% preamble, unchanged, then the scheme's own layer. A scheme's .tex is
% rendered from its YAML (scripts/README.md, Mark schemes), never edited.
\input{\CurrentFilePath/pack_preamble.tex}

% The page: equal side margins, so the table runs the full width of the
% header rule. The footers' right edge follows the text's.
\geometry{left=0.8in, right=0.8in}
\makeatletter
\renewcommand{\djs@footR}[1]{%
  \put(\dimexpr1in+\hoffset+\oddsidemargin+\textwidth\relax,
    -\dimexpr\paperheight-0.32in\relax){\makebox[0pt][r]{#1}}}
\makeatother

\usepackage[export]{adjustbox}
\usepackage{tabularray}
\definecolor{msGrid}{RGB}{176, 172, 164}
\definecolor{msTint}{RGB}{245, 243, 239}

% The header on every page, the front page included, and the booklet label
% under the contents-page title. \djsQuestionPageStyle stays empty.
\newcommand{\djsSchemeHeader}[1]{%
  \djsPackHeader[markscheme]{\textit{Mark scheme} - #1}%
  \renewcommand{\djsPackTitle}{#1}%
  \renewcommand{\djsBookletLabel}{Mark scheme}%
}

% The front page: the contents heading, spaced as tocloft spaces the packs'
% "Questions", then the contents table. A question's name and page number
% both link to its first page.
\newcommand{\msContentsHeading}{%
  \vskip\cftbeforetoctitleskip
  {\noindent\cfttoctitlefont\contentsname\par}%
  \vskip\cftaftertoctitleskip}
\NewTblrEnviron{mscontents}
\SetTblrInner[mscontents]{
  colspec = {Q[wd=26mm, l] X[l] Q[wd=16mm, c] Q[wd=14mm, r]},
  width = \linewidth,
  rowsep = 4pt, colsep = 6pt,
  hlines = {0.4pt, msGrid}, vlines = {0.4pt, msGrid},
  hline{2} = {2.2pt, scarlet},
  row{1} = {font=\small\bfseries},
  row{Z} = {font=\bfseries},
  column{2} = {font=\small},
}
\newcommand{\msContentsQuestion}[1]{\hyperref[q:#1]{Question #1}}
\newcommand{\msContentsPage}[1]{\hyperref[q:#1]{\pageref*{q:#1}}}
\newcommand{\msSep}{\hspace{0.9em}}

% The conventions pages' layout: headings in the documents' own bold, tables
% ruled and padded as the scheme's, roomy lists.
\newcommand{\msHead}[1]{%
  \par\vspace{18pt}%
  {\noindent\bfseries #1\par}%
  \vspace{6pt}}
\NewTblrEnviron{mstblr}
\SetTblrInner[mstblr]{rowsep = 6pt, colsep = 6pt,
  hlines = {0.4pt, msGrid}, vlines = {0.4pt, msGrid}}
\newenvironment{msList}
  {\begin{itemize}[leftmargin=1.2em, itemsep=5pt, topsep=3pt, parsep=0pt]}
  {\end{itemize}}

% A question's first page: the link target and the PDF bookmark, as
% \questionitem gives the packs.
\newcommand{\msQnum}{}
\newcommand{\msQuestion}[1]{%
  \clearpage
  \renewcommand{\msQnum}{#1}%
  \phantomsection\label{q:#1}%
  \addcontentsline{toc}{subsection}{\protect\djsTocEntry{#1}}}

% The scheme's table: Question 16mm, Scheme the rest, Marks 16mm, Notes
% 44mm; every cell ruled in msGrid; the scarlet rule under the head row. A
% long table has no caption or "continued" foot; a continuation page opens
% with a muted line above the repeated head row. The rules across a part's
% rows, the alternative's tint and the total row are the rendered table's.
\NewTblrTheme{ms}{%
  \DefTblrTemplate{firsthead}{default}{}%
  \DefTblrTemplate{middlehead,lasthead}{default}{%
    {\small\itshape\color{djsMuted}Question \msQnum\ continued\par}%
    \vspace{6pt}}%
  \DefTblrTemplate{firstfoot,middlefoot,lastfoot}{default}{}%
}
\NewTblrEnviron{msscheme}
\SetTblrOuter[msscheme]{long, theme=ms, entry=none, label=none}
\SetTblrInner[msscheme]{
  colspec = {Q[wd=16mm, l] X[l] Q[wd=16mm, c] Q[wd=44mm, l]},
  width = \linewidth,
  rowsep = 6pt, colsep = 6pt,
  rowhead = 1,
  vlines = {0.4pt, msGrid},
  hline{1} = {0.4pt, msGrid}, hline{2} = {2.2pt, scarlet},
  hline{Z} = {0.4pt, msGrid},
  column{1} = {font=\bfseries}, column{3} = {font=\bfseries},
  column{4} = {font=\small},
  row{1} = {rowsep=4pt, font=\small\bfseries},
}

% Two marks on a row: each note starts with its code.
\newcommand{\msNote}[2]{\textbf{#1}\enspace #2}
% A full-row figure: the worked solution's own picture, scaled uniformly and
% only down, to the table less the Question column. The picture is drawn
% with the line width it has in the solutions document, at its list depth
% there (the optional argument, 1 by default: a part's working), so one
% sized by \linewidth keeps its size; tabularray sets the cell's \linewidth
% to the cell's measured width, which would shrink it at every pass. The
% widths are measured once, in a solution box at the packs' text width
% (pack_preamble.tex's margins, 0.8in and 1in).
\newlength{\msSpanW}
\makeatletter
\newcommand{\ms@keepLW}[1]{%
  \expandafter\xdef\csname ms@solnLW@#1\endcsname{\the\linewidth}}
\AtBeginDocument{%
  \setlength{\msSpanW}{\dimexpr\textwidth-16mm-28pt\relax}%
  \setbox\z@\vbox{%
    \begin{minipage}{\dimexpr\paperwidth-1.8in\relax}
      \begin{djsSolution}
        \ms@keepLW{0}%
        \begin{enumerate}\item \ms@keepLW{1}%
        \begin{enumerate}\item \ms@keepLW{2}%
        \begin{enumerate}\item \ms@keepLW{3}%
        \end{enumerate}\end{enumerate}\end{enumerate}
      \end{djsSolution}
    \end{minipage}}}
\newcommand{\msFigure}[2][1]{%
  \adjustbox{max width=\msSpanW, max totalheight=0.6\textheight, valign=t}{%
    \hsize=\csname ms@solnLW@#1\endcsname
    \linewidth=\hsize
    #2}}
\makeatother
% A figure mark's feature, in the Scheme column.
\newcommand{\msFeature}[1]{\textit{#1}}
% An alternative's title row.
\newcommand{\msAlternative}[1]{\textit{Alternative:} #1}
% Drafts only: a queried mark's code, the query's why under its note, and a
% flag in its part's place.
\newcommand{\msQueried}[2]{#1~\(\to\) #2?}
\newcommand{\msQuery}[1]{{\itshape\color{djsMuted}Query: #1}}
\newcommand{\msFlag}[2]{\textbf{#1:}\enspace #2}


\djsSchemeHeader{Invariant Points and Lines}
\djsSitePath{Core Pure}{Vectors, Matrices \& Linear Transformations}

\djsLastUpdated{25 October 2026}
\begin{document}
\djsSetQuestionCount{9}
\djsAddMarks{1}{4}
\djsAddMarks{2}{5}
\djsAddMarks{3}{9}
\djsAddMarks{4}{9}
\djsAddMarks{5}{13}
\djsAddMarks{6}{8}
\djsAddMarks{7}{7}
\djsAddMarks{8}{7}
\djsAddMarks{9}{13}
\vspace*{6pt}
\djsContentsTitle
\vspace{30pt}
\msContentsHeading
\begin{mscontents}{}
Question & Marks by part & Marks & Page \\
\msContentsQuestion{1} & (a)\,2\msSep (b)\,2 & 4 & \msContentsPage{1} \\
\msContentsQuestion{2} & & 5 & \msContentsPage{2} \\
\msContentsQuestion{3} & (a)\,4\msSep (b)\,5 & 9 & \msContentsPage{3} \\
\msContentsQuestion{4} & (a)\,2\msSep (b)\,2\msSep (c)\,3\msSep (d)\,2 & 9 & \msContentsPage{4} \\
\msContentsQuestion{5} & (a)\,4\msSep (b)\,2\msSep (c)\,7 & 13 & \msContentsPage{5} \\
\msContentsQuestion{6} & (a)\,5\msSep (b)\,3 & 8 & \msContentsPage{6} \\
\msContentsQuestion{7} & (a)\,3\msSep (b)\,4 & 7 & \msContentsPage{7} \\
\msContentsQuestion{8} & (a)\,5\msSep (b)\,2 & 7 & \msContentsPage{8} \\
\msContentsQuestion{9} & (i)\,8\msSep (ii)\,2\msSep (iii)\,3 & 13 & \msContentsPage{9} \\
\SetCell[c=2]{r} Total & & 75 & \\
\end{mscontents}
\newpage
% The conventions pages, shared by every pack's mark scheme: one source of
% truth, \input after the front page, with no strand or pack logic. The keys
% on the first side, the rules on the second. The scheme layer defines
% \msHead, msList and mstblr.
\vspace*{6pt}%
{\noindent\cfttoctitlefont Reading the mark scheme\par}%
\vspace{10pt}%
These rules apply throughout, unless a note says otherwise.

\msHead{Marks}
Every mark is worth one, so \textbf{M1 A1} on a row is two marks. Only marks
coded ft follow through.\par
\vspace{10pt}%
\begin{mstblr}{
  colspec = {Q[wd=15mm, l] X[l]},
  width = \linewidth,
  hline{2} = {2.2pt, scarlet},
  column{1} = {font=\bfseries},
  row{1} = {font=\small\bfseries, rowsep=4pt},
}
Code & Meaning \\
M1 & Method: a correct method, applied to this question. A slip doesn't cost it if the method is clear. A formula must be used, not just quoted. \\
dM1 & Dependent method: as M1, but only if the method mark before it was earned, unless its note names another. \\
A1 & Accuracy: a correct answer or step. It needs the method mark it follows. \\
A1ft & As A1, following through from an earlier answer that the note names. \\
A1~AG & The last mark of a part whose answer is printed in the question: complete working, with no errors. \\
B1 & A correct value, fact or statement. It needs no method mark. \\
B1ft & As B1, following through from an earlier answer that the note names. \\
E1 & A written explanation, interpretation or conclusion. Concluding a proof by induction needs every earlier mark. \\
\end{mstblr}

\msHead{Abbreviations}
ft and AG qualify a code in the Marks column. The others appear in the
notes.\par
\vspace{10pt}%
\begin{mstblr}{
  colspec = {Q[wd=17mm, l] X[l] Q[wd=15mm, l] X[l]},
  width = \linewidth,
  column{1} = {font=\bfseries}, column{3} = {font=\bfseries},
}
ft & follow through & AG & answer given in the question \\
cao & correct answer only, with no follow-through & cso & correct solution only: no errors in the part \\
oe & or an equivalent form & awrt & any answer that rounds to this value \\
soi & seen, or implied by later working & isw & ignore working after a correct answer \\
SC & special case: credit outside the main method & their & the student's own earlier value \\
condone & a slip that doesn't cost the mark & [\;\,] & working the mark doesn't need \\
s.f. & significant figures & d.p. & decimal places \\
\end{mstblr}

\newpage
\msHead{Reading the scheme}
\begin{msList}
\item Each row is a line of working, with the marks it earns beside it.
  Where a row carries two marks, each note starts with its code.
\item \textbf{[\textit{n}]} ends each part with its total.
\item A shaded \textit{Alternative} is a second method for the same total.
  Mark an answer by the method it follows.
\item Any other valid method earns the equivalent marks, unless the question
  names a method or says \textit{hence}.
\item A sketch earns one B1 for each feature named below its figure.
\end{msList}

\msHead{Accuracy and answers}
\begin{msList}
\item Unless a question says otherwise, give answers to 3 significant
  figures, and angles in degrees to 1 decimal place. Angles in radians take
  3 significant figures.
\item A final answer with more figures is accepted if it rounds to the
  scheme's value; one with fewer is not. An answer spoiled by rounding too
  early loses its accuracy mark.
\item Where a question asks for an exact answer, a decimal loses the final
  accuracy mark. Otherwise either is accepted.
\item Wrong or missing units cost nothing unless the question asks for them.
\item A correct answer with no working earns full marks, unless the question
  needs the working: show that, prove, hence, a method it names, or an exact
  answer. A wrong answer with no working earns nothing. Calculator answers
  are treated the same way.
\item A correct answer, once reached, isn't spoiled by later working (isw).
\end{msList}

\msHead{Value of $\boldsymbol{g}$}
\begin{msList}
\item Take $g = 9.8\,\mathrm{m\,s^{-2}}$ unless a question states a value,
  which is then the only one accepted. Answers that use $g = 9.8$ may be
  given to 2 or 3 significant figures.
\item Where a question doesn't state $g$, an answer worked with $g = 10$
  throughout earns the same marks, keeping any value printed in the
  question. The worked solutions list the final answers with $g = 10$.
\end{msList}


\msQuestion{1}
\begin{msscheme}{
  hline{3-4,6-7} = {2-Z}{0.4pt, msGrid},
  hline{5,8} = {0.4pt, msGrid},
  row{8} = {font=\bfseries},
}
Question & Scheme & Marks & Notes \\
1(a) & {\(p = -1\): \(2(3) - 2(4) = \ldots\), \(1(3) + 2(4) = \ldots\)} & M1 & {Substitutes \(p = -1\) into \(\mathbf{M}\) and multiplies by the column vector \((3, 4)\). Condone one slip in an entry.} \\
& {\(6 - 8 = -2\), \(3 + 8 = 11 \Rightarrow (-2, 11)\)} & A1~AG & {Correct matrix and both products leading without error to \((-2, 11)\).} \\
\nopagebreak & & [2] & \\
1(b) & {\(p = 2\): \(5x - 2y = x\), \(4x - y = y\)} & M1 & {Substitutes \(p = 2\) and sets \(\mathbf{M}\binom{x}{y} = \binom{x}{y}\), forming at least one equation in \(x\) and \(y\). Condone one slip in an entry.} \\
& {\(4x - 2y = 0 \Rightarrow y = 2x\), gradient \(2\)} & A1 & \\
\nopagebreak & & [2] & \\
\nopagebreak \SetCell[c=2]{r} Total for Question 1 & & 4 & \\
\end{msscheme}

\msQuestion{2}
\begin{msscheme}{
  hline{3-5} = {2-Z}{0.4pt, msGrid},
  hline{6} = {0.4pt, msGrid},
  row{6} = {font=\bfseries},
}
Question & Scheme & Marks & Notes \\
2 & {\((x,\ mx+c) \mapsto \bigl((1+4m)x+4c,\ (2+3m)x+3c\bigr)\)\\[4pt] \((2+3m)x+3c = m\bigl((1+4m)x+4c\bigr)+c\)} & M1 & {Substitutes the image of a general point of \(y=mx+c\) into \(y=mx+c\). Condone slips in the matrix product.} \\
& {\(2+3m = m+4m^2\), \(\ 3c = (4m+1)c\)} & M1 A1 & {\msNote{M1}{Equates coefficients of \(x\) and constant terms in their equation.}\\ \msNote{A1}{Both equations correct, oe.}} \\
& {\(2m^2-m-1=0 \Rightarrow m = \ldots\)} & M1 & {Solves their three-term quadratic in \(m\) by factorising, formula, completing the square or calculator.} \\
& {\(m=-\tfrac12\), \(c=0 \Rightarrow y=-\tfrac12x\)} & A1 & {\(y=-\frac12x\) oe, with \(c=0\) found from the constant equation.} \\
\nopagebreak \SetCell[c=2]{r} Total for Question 2 & & 5 & \\
\end{msscheme}

\msQuestion{3}
\begin{msscheme}{
  hline{3-6,8-12} = {2-Z}{0.4pt, msGrid},
  hline{7,13} = {0.4pt, msGrid},
  row{13} = {font=\bfseries},
}
Question & Scheme & Marks & Notes \\
3(a) & {\((p+2)(p-1) - (-5)(-2) = 0\)} & M1 & {Forms \(ad - bc\) for \(\mathbf{M}\) in terms of \(p\) and sets it equal to zero. Condone one sign slip.} \\
& {\(p^2 + p - 12 = 0\)} & A1 & {Correct three-term quadratic oe.} \\
& {\((p+4)(p-3) = 0 \Rightarrow p = \ldots\)} & dM1 & {Solves their three-term quadratic by factorising with constants multiplying to their constant term, the formula or completing the square; calculator allowed.} \\
& {\(p = 3\) or \(p = -4\)} & A1 & {Both values, and no others.} \\
\nopagebreak & & [4] & \\
3(b) & {\(\begin{pmatrix} 5 & -5 \\ -2 & 2 \end{pmatrix}\begin{pmatrix} x \\ mx \end{pmatrix} = \begin{pmatrix} 5x - 5mx \\ -2x + 2mx \end{pmatrix}\)} & M1 & {Applies \(\mathbf{M}\) to a general point of \(y = mx\) and requires the image to satisfy \(y = mx\). Condone slips in the product.} \\
& {\(-2x + 2mx = m(5x - 5mx)\)} & A1 & {Correct equation oe; \(x\) may already be cancelled.} \\
& {\(5m^2 - 3m - 2 = 0 \Rightarrow m = \ldots\)} & dM1 & {Rearranges, after cancelling \(x\), to a three-term quadratic in \(m\) and solves it by factorising, the formula or completing the square; calculator allowed.} \\
& {\(y = x\)} & A1 & {An equation of the line; \(m = 1\) alone is not enough.} \\
& {\(y = -\frac25 x\)} & A1 & {oe, such as \(2x + 5y = 0\); \(m = -\frac25\) alone is not enough.} \\
\nopagebreak & & [5] & \\
\nopagebreak \SetCell[c=2]{r} Total for Question 3 & & 9 & \\
\end{msscheme}

\msQuestion{4}
\begin{msscheme}{
  hline{3-4,6-7,9-11,13-14} = {2-Z}{0.4pt, msGrid},
  hline{5,8,12,15} = {0.4pt, msGrid},
  row{15} = {font=\bfseries},
}
Question & Scheme & Marks & Notes \\
4(a) & {A reflection leaves every point on its mirror line unchanged} & E1 & {States that points on the mirror line are unchanged by the reflection. Not for asserting only that a reflection has invariant points.} \\
& {The line of invariant points is the mirror line} & E1 & {States that it is the mirror line; accept axis or line of reflection.} \\
\nopagebreak & & [2] & \\
4(b) & {\(\tfrac15(-3x - 4y) = x,\quad \tfrac15(-4x + 3y) = y\)\\[4pt] \(\Rightarrow 2x + y = 0\)} & M1 & {Equates the image of \((x, y)\) under \(\mathbf{A}\) to \((x, y)\) and forms at least one linear equation in \(x\) and \(y\). Condone sign slips.} \\
& {\(y = -2x\)} & A1 & {Accept \(y = -2x + 0\). A0 for \(2x + y = 0\) alone.} \\
\nopagebreak & & [2] & \\
4(c) & {\(P' = \tfrac15(-12 - 28,\ -16 + 21)\ [= (-8, 1)]\)} & M1 & {Multiplies \(\mathbf{A}\) by the position vector of \(P\). Condone one arithmetic slip.} \\
& {\(\left(\frac{4 + \their{{-8}}}{2},\ \frac{7 + \their{1}}{2}\right)\)} & dM1 & {Finds the midpoint of \(P\) and their image.} \\
& {\((-2, 4)\)} & A1 & \\
\nopagebreak & & [3] & \\
4(d) & {Perpendicular to the mirror line:\\[4pt] \((-2)\left(-\tfrac2b\right) = -1 \Rightarrow b = \ldots\)} & M1 & {Sets the product of the gradient of the line and their mirror-line gradient equal to \(-1\), oe. Condone a sign slip in \(-\frac2b\).} \\
& {\(b = -4\)} & A1 & \\
\nopagebreak & & [2] & \\
\nopagebreak \SetCell[c=2]{r} Total for Question 4 & & 9 & \\
\end{msscheme}

\msQuestion{5}
\begin{msscheme}{
  hline{3-5,7-8,10-15} = {2-Z}{0.4pt, msGrid},
  hline{6,9,16} = {0.4pt, msGrid},
  row{16} = {font=\bfseries},
}
Question & Scheme & Marks & Notes \\
5(a) & {Reflection in \(y = \frac{x}{\sqrt3}\)} & B1 B1 & {\msNote{B1}{Reflection.}\\ \msNote{B1}{In \(y = \frac{x}{\sqrt3}\); accept \(y = \frac{\sqrt3}{3}x\), \(y = x\tan 30^\circ\), or the line through the origin at \(30^\circ\) or \(\frac{\pi}{6}\) to the \(x\)-axis.}} \\
& {Stretch parallel to the \(y\)-axis, scale factor 6} & B1 & {Accept in the \(y\)-direction, or with the \(x\)-axis invariant.} \\
& {Reflection first, then the stretch} & B1 & {Order stated explicitly, with both transformations named.} \\
\nopagebreak & & [4] & \\
5(b) & {Inverts each factor, in reverse order:\\[4pt] \(\mathbf{M}^{-1} = \begin{pmatrix}\frac12 & \frac{\sqrt3}{2}\\ \frac{\sqrt3}{2} & -\frac12\end{pmatrix}^{-1}\begin{pmatrix}1 & 0\\ 0 & 6\end{pmatrix}^{-1}\)} & M1 & {Multiplies the inverses of the two factors in reverse order; condone slips in either inverse. Not for the original order.} \\
& {\(\mathbf{M}^{-1} = \begin{pmatrix}\frac12 & \frac{\sqrt3}{2}\\ \frac{\sqrt3}{2} & -\frac12\end{pmatrix}\begin{pmatrix}1 & 0\\ 0 & \frac16\end{pmatrix}\)} & A1 & {Both inverses evaluated, in this order. A0 for a single multiplied-out matrix.} \\
\nopagebreak & & [2] & \\
5(c) & {\(\mathbf{M} = \begin{pmatrix}\frac12 & \frac{\sqrt3}{2}\\ 3\sqrt3 & -3\end{pmatrix}\)} & B1 & \\
& {\(\their{\mathbf{M}}\begin{pmatrix}1\\ m\end{pmatrix} = k\begin{pmatrix}1\\ m\end{pmatrix}\)\\[4pt] \(\Rightarrow \frac12 + \frac{\sqrt3}{2}m = k,\ 3\sqrt3 - 3m = km\)} & M1 & {Applies their \(\mathbf{M}\) to \((x, mx)\) or \((1, m)\) and sets the image on \(y = mx\), giving two equations.} \\
& {\(3\sqrt3 - 3m = m\left(\frac12 + \frac{\sqrt3}{2}m\right)\)} & M1 & {Eliminates \(k\), or substitutes the image into \(y = mx\), to form an equation in \(m\) alone.} \\
& {\(\sqrt3 m^2 + 7m - 6\sqrt3 = 0\)} & A1 & {oe} \\
& {\((\sqrt3 m - 2)(m + 3\sqrt3) = 0 \Rightarrow m = \ldots\)} & M1 & {Solves their three-term quadratic by factors whose constants multiply to their constant term, the formula with their values, or completing the square. Calculator allowed.} \\
& {\(y = \frac{2}{\sqrt3}x,\quad y = -3\sqrt3 x\)} & A1 A1 & {\msNote{A1}{\(y = \frac{2}{\sqrt3}x\) oe, or awrt \(y = 1.15x\). An equation, not \(m\) alone.}\\ \msNote{A1}{\(y = -3\sqrt3 x\) oe, or awrt \(y = -5.20x\). An equation, not \(m\) alone.}} \\
\nopagebreak & & [7] & \\
\nopagebreak \SetCell[c=2]{r} Total for Question 5 & & 13 & \\
\end{msscheme}

\msQuestion{6}
\begin{msscheme}{
  hline{3-7,9-11} = {2-Z}{0.4pt, msGrid},
  hline{8,12} = {0.4pt, msGrid},
  row{12} = {font=\bfseries},
}
Question & Scheme & Marks & Notes \\
6(a) & {\(\begin{pmatrix}2&-1\\ p&0\end{pmatrix}\begin{pmatrix}0\\y\end{pmatrix} = \begin{pmatrix}-y\\0\end{pmatrix}\)\\[4pt] not on \(x = 0\) for \(y \neq 0\), so not invariant} & B1 & {Image of a general point, or of a specific point other than the origin, with a statement that it is not on \(x = 0\).} \\
& {\(\begin{pmatrix}2&-1\\ p&0\end{pmatrix}\begin{pmatrix}x\\mx\end{pmatrix} = \begin{pmatrix}(2-m)x\\px\end{pmatrix}\)\\[4pt] \(\Rightarrow px = m(2-m)x\)} & M1 & {Finds the image of a general point \((x, mx)\) and substitutes it into \(y = mx\). Condone sign slips in the image.} \\
& {\(\Rightarrow m^2 - 2m + p = 0\)} & A1 & {Correct substituted equation leading without error to the printed result; combined steps accepted.} \\
& {No invariant lines \(\Leftrightarrow (-2)^2 - 4p < 0\)} & M1 & {Sets the discriminant of the quadratic in \(m\) negative, or completes the square. Condone slips in the discriminant.} \\
& {\(p > 1\)} & A1 & {Strict inequality, oe in set notation.} \\
\nopagebreak & & [5] & \\
6(b) & {\(\det\mathbf{M} = p = -3\)} & B1 & \\
& {\(m^2 - 2m - 3 = 0 \Rightarrow m = \ldots\)} & M1 & {Substitutes their value of \(p\) into \(m^2 - 2m + p = 0\) and solves the quadratic by any means; a calculator is allowed.} \\
& {\(y = 3x\) and \(y = -x\)} & A1 & {Both, as equations of lines. Lines not through the origin need not be considered.} \\
\nopagebreak & & [3] & \\
\nopagebreak \SetCell[c=2]{r} Total for Question 6 & & 8 & \\
\end{msscheme}

\msQuestion{7}
\begin{msscheme}{
  hline{3-5,7-10} = {2-Z}{0.4pt, msGrid},
  hline{6,11} = {0.4pt, msGrid},
  row{11} = {font=\bfseries},
}
Question & Scheme & Marks & Notes \\
7(a) & {\((x, mx) \mapsto \big(x - 4mx,\ 2x - mx\big)\)} & M1 & {Applies \(\mathbf{B}\) to a general point of the line, such as \((x, mx)\) or \((1, m)\), to find its image. Condone one slip.} \\
& {\(\dfrac{Y}{X} = \dfrac{(2-m)x}{(1-4m)x} \Rightarrow y = \dfrac{2-m}{1-4m}x\)} & A1~AG & {Correct image coordinates leading without error to the printed equation, by their ratio or by eliminating \(x\).} \\
& {\(y = \tfrac14x \mapsto x = 0\)\\[4pt] \(x = 0 \mapsto y = \tfrac14x\)} & B1 & {Both images, oe such as \(4y = x\) or the \(y\)-axis.} \\
\nopagebreak & & [3] & \\
7(b) & {\(m = \dfrac{2-m}{1-4m}\)} & M1 & {Equates \(m\) to the image gradient from (a).} \\
& {\(4m^2 - 2m + 2 = 0\)} & A1 & {oe, such as \(2m^2 - m + 1 = 0\).} \\
& {\((-2)^2 - 4(4)(2) = \ldots\)} & M1 & {Evaluates the discriminant of their three-term quadratic, or completes the square, to test for real roots.} \\
& {\(-28 < 0 \Rightarrow\) no real \(m\)\\[4pt] \(y = \tfrac14x\) and \(x = 0\) map to each other,\\[4pt] so neither is invariant\\[4pt] No invariant line through the origin, so none\\[4pt] parallel to one: \(T\) has no invariant lines} & A1~AG & {Correct negative discriminant, both special lines shown not invariant from (a), and a conclusion citing the parallel-line fact. A0 if either special line is omitted.} \\
\nopagebreak & & [4] & \\
\nopagebreak \SetCell[c=2]{r} Total for Question 7 & & 7 & \\
\end{msscheme}

\msQuestion{8}
\begin{msscheme}{
  hline{3-7,9-10} = {2-Z}{0.4pt, msGrid},
  hline{8,11} = {0.4pt, msGrid},
  row{11} = {font=\bfseries},
}
Question & Scheme & Marks & Notes \\
8(a) & {\(\mathbf{A}^2 = \begin{pmatrix} b^2+(a-3)(2a-1) & (a-3)(b-2) \\ (2a-1)(b-2) & (2a-1)(a-3)+4 \end{pmatrix}\)} & M1 & {Multiplies \(\mathbf{A}\) by itself and sets the result equal to \(\mathbf{I}\), or equates \(\mathbf{A}\) to its inverse by formula. Condone slips in entries.} \\
& {\(b^2+(a-3)(2a-1) = 1\), \((a-3)(b-2) = 0\)\\[4pt] \((2a-1)(b-2) = 0\), \((2a-1)(a-3)+4 = 1\)} & A1 & {All four equations correct; need not be simplified.} \\
& {\(2a^2 - 7a + 6 = 0 \Rightarrow a = \ldots\)} & dM1 & {Forms a three-term quadratic in \(a\) from the bottom-right entry and solves it by any means.} \\
& {\(a = \frac32\) or \(a = 2\)} & A1 & {Both values, and no others.} \\
& {\(a \neq 3 \Rightarrow (a-3)(b-2) = 0 \Rightarrow b = 2\)} & A1 & {\(b = 2\) only, from an off-diagonal equation, with \(a - 3\) or \(2a - 1\) shown or stated non-zero for both values of \(a\).} \\
\nopagebreak & & [5] & \\
8(b) & {\(\begin{pmatrix} 2 & -\frac32 \\ 2 & -2 \end{pmatrix}\begin{pmatrix} x \\ y \end{pmatrix} = \begin{pmatrix} x \\ y \end{pmatrix}\)} & M1 & {Uses \(\mathbf{A}\) with their smaller \(a\) and their \(b\) from (a), sets the image of \((x, y)\) equal to \((x, y)\) and forms two equations.} \\
& {\(2x - \frac32 y = x \Rightarrow 2x - 3y = 0\)\\[4pt] \(2x - 2y = y \Rightarrow 2x - 3y = 0\)\\[4pt] Same equation, so a line: \(2x - 3y = 0\)} & A1 & {Both equations reduced to the same line and stated equivalent, giving \(2x - 3y = 0\) oe, such as \(y = \frac23 x\).} \\
\nopagebreak & & [2] & \\
\nopagebreak \SetCell[c=2]{r} Total for Question 8 & & 7 & \\
\end{msscheme}

\msQuestion{9}
\begin{msscheme}{
  hline{3-10,12-13,15-17} = {2-Z}{0.4pt, msGrid},
  hline{11,14,18} = {0.4pt, msGrid},
  row{18} = {font=\bfseries},
}
Question & Scheme & Marks & Notes \\
9(i) & {\(\begin{pmatrix}a&b\\c&d\end{pmatrix}\begin{pmatrix}1\\-1\end{pmatrix} = \begin{pmatrix}2\\0\end{pmatrix}\)} & M1 & {Equates the image of \((1,-1)\) under a general matrix to \((2,0)\), forming equations in the entries. Condone one slip.} \\
& {\(a - b = 2,\quad c - d = 0\)} & A1 & {Both, oe.} \\
& {\(|\det\mathbf A| = \frac{24}{4} \Rightarrow |ad - bc| = 6\)} & M1 & {Equates \(\det\mathbf A\) or \(|\det\mathbf A|\) to \(\frac{24}{4}\); \(\det\mathbf A = 6\) alone, or \(-6\) alone, is enough.} \\
& {\((b + 2)d - bd = \pm 6 \Rightarrow 2d = \pm 6\)} & A1 & {Correct equation in the remaining entries, oe such as \(2c = \pm 6\); one sign alone accepted.} \\
& {\(\det(\mathbf A - \mathbf I) = (a - 1)(d - 1) - bc = 0\)} & M1 & {Uses \(\det(\mathbf A - \mathbf I) = 0\), or another valid condition for a line of invariant points, with their entries. Condone sign slips.} \\
& {\(d - b - 1 = 0\)} & A1 & {oe, in any entries.} \\
& {\(d = 3 \Rightarrow b = \ldots,\ a = \ldots,\ c = \ldots\)} & dM1 & {Solves their three equations together, choosing a sign where needed, to find all four entries.} \\
& {\(\mathbf A = \begin{pmatrix}4&2\\3&3\end{pmatrix}\)} & A1 & {Or \(\begin{pmatrix}-2&-4\\-3&-3\end{pmatrix}\).} \\
\nopagebreak & & [8] & \\
9(ii) & {\(\begin{pmatrix}2&1\\2&3\end{pmatrix}\begin{pmatrix}x\\y\end{pmatrix} = \begin{pmatrix}x\\y\end{pmatrix} \Rightarrow 2x + y = x\)} & M1 & {Sets the image of \((x, y)\) under \(\mathbf B\) equal to \((x, y)\) and forms at least one equation. Condone one slip.} \\
& {\(y = -x\)} & A1 & {oe such as \(x + y = 0\).} \\
\nopagebreak & & [2] & \\
9(iii) & {\(\begin{pmatrix}2&1\\2&3\end{pmatrix}\begin{pmatrix}x\\2x + c\end{pmatrix} = \ldots\)} & M1 & {Multiplies \(\mathbf B\) by a general point of the line, such as \((t, 2t + c)\). Condone one slip.} \\
& {\(\phantom{A} = \begin{pmatrix}4x + c\\8x + 3c\end{pmatrix}\)} & A1 & \\
& {\(2(4x + c) + c = 8x + 3c\)\\[4pt] so the image lies on \(y = 2x + c\): invariant} & A1~AG & {Correct substitution showing the image satisfies \(y = 2x + c\), without error, and a conclusion that the line is invariant.} \\
\nopagebreak & & [3] & \\
\nopagebreak \SetCell[c=2]{r} Total for Question 9 & & 13 & \\
\end{msscheme}

\end{document}
